% Block 0 of the slides (00_front.tex).
\section{The thesis: verification is the bottleneck}\label{sec:thesis}

\onslides{two motivation slides (TheoPhysEval, Tao), the thesis, and the
``new PhD student'' mental model.}

\subsection{Two data points from October 2026}

\textbf{Agents solve research problems.} Ramesh Chandra et al.\ \rf{theophyseval}
collected 50 unpublished research-level problems in theoretical physics, each with a
verifiable final answer and each about a day's work or less for a world expert. Agents
received literature search and the ability to write and run scripts, and no human hints.
The best agent reached $98\,\%$ correct final answers (49 of 50) at an average cost under
3\,USD per problem; the runner-up solved 46. Accuracy changed little when web search
and background papers were withheld, which the authors read as substantial memorisation
of the literature. They also discuss why no agent-led breakthrough has appeared despite
this problem-solving ability. Two points matter for the course: the problems had answers
that could be \emph{checked}, and someone had to build that check.

\textbf{Abundance outpaces verification.} Tao's essay \rf{tao}, based on his public
lecture at the 2026 International Congress of Mathematicians, assumes that AI tools will
perform research-level mathematical tasks and asks what mathematicians actually value.
In Section~7 he lists the ``impedance mismatches'' that would follow; the first is the
sentence quoted on the slide, that AI-generated proofs will accumulate faster than they
can be verified. He calls the result a transition from proof scarcity to proof
abundance, and notes that journals, priority conventions, hiring and prizes were all
designed for scarcity. The slide carries the argument to physics, where code, numerics
and derivations have become cheap in the same way.

\subsection{From the data to the thesis}

If producing a result costs minutes, the scarce resource is the confidence that the
result is right. Research software engineering (RSE) is the body of practice that builds
that confidence for code written by humans: version control, a fixed project structure,
automated tests and continuous integration \rf{wilson-good,wilson-best}. None of it was
designed for AI, and all of it applies, because every one of these practices was
invented for contributors who make mistakes. An agent is such a contributor, only
faster.

\subsection{The new PhD student}

The mental model on the slide has four traits, and each maps onto a later block.
\begin{itemize}
  \item \textbf{Fast}: it produces more than you can read, so the reading has to be
        made efficient (diffs, Section~\ref{sec:git}; tests, Section~\ref{sec:verification}).
  \item \textbf{Eager}: it rarely declines a task, and rarely says ``I don't know''
        (Section~\ref{sec:llm}, post-training).
  \item \textbf{Overconfident}: correct and incorrect output read the same
        (Section~\ref{sec:llm}, plausibility).
  \item \textbf{Forgetful}: nothing survives the session unless it is written to disk
        (Section~\ref{sec:context}).
\end{itemize}
The analogy has a limit worth stating if asked: a PhD student learns from your
corrections and becomes more reliable over months; a model does not, unless the
correction is written into a file it reads every time.

\subsection{A sceptical counterweight}

A randomised controlled trial by METR \rf{metr} gave 16 experienced open-source
developers 246 tasks in their own repositories, with AI tools allowed or disallowed per
task. Before the study they forecast that AI would cut completion time by $24\,\%$;
afterwards they estimated it had cut it by $20\,\%$; measured, it \emph{increased}
completion time by $19\,\%$. The study used early-2025 tools on large, mature codebases
the developers knew well, and the authors caution against generalising. Its lasting
lesson is the gap between perceived and measured productivity: how fast the work
\emph{feels} is not evidence. The slides dropped this study for pace; it is a good answer
to ``but does it really make us faster?''

\begin{qa}
\Q{Is TheoPhysEval not just measuring memorisation?}
Partly, and the authors say so: removing web search and papers changed accuracy little.
The problems were unpublished, but their ingredients were not. The course's point
survives either way: the answers were checkable, and the check made the result usable.

\Q{If agents are this good, why spend an afternoon on verification?}
Because ``98\,\%'' was measured against known answers. In your research there is no
answer key; the tests and cross-checks you write are the answer key.

\Q{Did Tao say AI is bad for mathematics?}
No. The essay conditions on capable AI arriving and asks which values and institutions
need to be made explicit. The quoted line describes a bottleneck, not a verdict.
\end{qa}

\textbf{Read next:} Tao \rf{tao}, Sections~1 and~7 (short and non-technical); the
TheoPhysEval abstract and discussion \rf{theophyseval}; the METR paper \rf{metr}.
